\documentclass[preview]{standalone}

\usepackage{amsmath}
\usepackage{amssymb}
\usepackage{stellar}
\usepackage{definitions}
\usepackage{bettelini}

\begin{document}

\id{probability-variance-covariance}
\genpage

\section{Moments}

\begin{snippetdefinition}{finite-moments-discrete-random-variable-definition}{Finite moments}
    Let \((\Omega,\powerset(\Omega),\mathbb P)\) be a
    \countableprobabilityspace, let
    \(X\colon\Omega\fromto E\subseteq\realnumbers\) be a real-valued
    \discreterandomvariable, and let \(k\in\naturalnumbers^\exceptzero\).
    The random variable \(X\) has a \emph{finite \(k\)-th moment} if
    \[
        \expectation[|X|^k]<+\infty.
    \]
    In this case, \(\expectation[X^k]\) is called the
    \emph{\(k\)-th moment} of \(X\).
\end{snippetdefinition}

\begin{snippetdefinition}{mixed-moment-discrete-random-variables-definition}{Mixed moment}
    Let \((\Omega,\powerset(\Omega),\mathbb P)\) be a
    \countableprobabilityspace, and let
    \(X,Y\colon\Omega\fromto\realnumbers\) be real-valued
    \discreterandomvariable[discrete random variables]. The pair \(X,Y\) has
    a \emph{finite mixed moment} if
    \[
        \expectation[|XY|]<+\infty.
    \]
    In this case, \(\expectation[XY]\) is called the
    \emph{mixed moment} of \(X\) and \(Y\).
\end{snippetdefinition}

\begin{snippetproposition}{discrete-moment-density-formulas}{Moment formulas through densities}
    Let \((\Omega,\powerset(\Omega),\mathbb P)\) be a
    \countableprobabilityspace, and let
    \(X,Y\colon\Omega\fromto\realnumbers\) be real-valued
    \discreterandomvariable[discrete random variables]. Let \(p_X\) be the
    \randomvariabledensity of \(X\), and let \(p_{X,Y}\) be the
    \jointdensity of \((X,Y)\). Then:
    \begin{enumerate}
        \item for \(k\in\naturalnumbers^\exceptzero\), \(X\) has a
            \finitemoment if and only if
            \[
                \sum_{x\in\realnumbers}|x|^kp_X(x)<+\infty,
            \]
            and in this case
            \[
                \expectation[X^k]
                =
                \sum_{x\in\realnumbers}x^kp_X(x);
            \]
        \item \(X,Y\) have a \mixedmoment if and only if
            \[
                \sum_{(x,y)\in\realnumbers^2}|xy|p_{X,Y}(x,y)<+\infty,
            \]
            and in this case
            \[
                \expectation[XY]
                =
                \sum_{(x,y)\in\realnumbers^2}xyp_{X,Y}(x,y).
            \]
    \end{enumerate}
\end{snippetproposition}

\begin{snippetproof}{discrete-moment-density-formulas-proof}{discrete-moment-density-formulas}{Moment formulas through densities}
    Apply the change-of-variable formula for the \expectedvalue to
    \(g(x)=x^k\) and to \(g(x,y)=xy\).
\end{snippetproof}

\begin{snippetproposition}{lower-moments-from-higher-moment}{Lower moments from a finite higher moment}
    Let \((\Omega,\powerset(\Omega),\mathbb P)\) be a
    \countableprobabilityspace, and let
    \(X\colon\Omega\fromto\realnumbers\) be a real-valued
    \discreterandomvariable. If
    \[
        \expectation[|X|^k]<+\infty
    \]
    for some \(k\in\naturalnumbers^\exceptzero\), then
    \[
        \expectation[|X|^h]<+\infty
    \]
    for every \(h\in\{1,\ldots,k\}\).
\end{snippetproposition}

\begin{snippetproof}{lower-moments-from-higher-moment-proof}{lower-moments-from-higher-moment}{Lower moments from a finite higher moment}
    For every \(x\in\realnumbers\) and every \(h\leq k\),
    \[
        |x|^h\leq 1+|x|^k.
    \]
    Taking expected values gives the result.
\end{snippetproof}

\section{Inequalities}

\begin{snippettheorem}{discrete-jensen-inequality}{Jensen's inequality}
    Let \((\Omega,\powerset(\Omega),\mathbb P)\) be a
    \countableprobabilityspace, let
    \(X\colon\Omega\fromto\realnumbers\) be a real-valued
    \discreterandomvariable with finite \expectedvalue, and let
    \(f\colon\realnumbers\fromto\realnumbers\) be a convex \function. If
    \(f(X)\) has an \expectedvalue, then
    \[
        f(\expectation[X])\leq\expectation[f(X)].
    \]
\end{snippettheorem}

\begin{snippetproof}{discrete-jensen-inequality-proof}{discrete-jensen-inequality}{Jensen's inequality}
    A convex real \function admits a supporting affine function at every
    point of the interior of its domain. Hence there exists
    \(\lambda\in\realnumbers\) such that, for every \(x\in\realnumbers\),
    \[
        f(x)\geq f(\expectation[X])+\lambda(x-\expectation[X]).
    \]
    Substituting \(X\) and taking expected values gives
    \[
        \expectation[f(X)]
        \geq
        f(\expectation[X])+\lambda(\expectation[X]-\expectation[X])
        =
        f(\expectation[X]).
    \]
\end{snippetproof}

\begin{snippettheorem}{probability-cauchy-schwarz-inequality}{Cauchy-Schwarz inequality for discrete random variables}
    Let \((\Omega,\powerset(\Omega),\mathbb P)\) be a
    \countableprobabilityspace, and let
    \(X,Y\colon\Omega\fromto\realnumbers\) be real-valued
    \discreterandomvariable[discrete random variables]. If
    \[
        \expectation[X^2]<+\infty,
        \qquad
        \expectation[Y^2]<+\infty,
    \]
    then \(X,Y\) have a \mixedmoment and
    \[
        |\expectation[XY]|
        \leq
        \sqrt{\expectation[X^2]\expectation[Y^2]}.
    \]
\end{snippettheorem}

\begin{snippetproof}{probability-cauchy-schwarz-inequality-proof}{probability-cauchy-schwarz-inequality}{Cauchy-Schwarz inequality for discrete random variables}
    Since \(2|XY|\leq X^2+Y^2\), the mixed moment is finite. If
    \(\expectation[X^2]=0\) or \(\expectation[Y^2]=0\), then one of the
    random variables is \(0\) almost surely, and the result is immediate.
    Otherwise, for every \(t\in\realnumbers\),
    \[
        0\leq\expectation[(X+tY)^2]
        =
        \expectation[X^2]+2t\expectation[XY]+t^2\expectation[Y^2].
    \]
    This quadratic polynomial in \(t\) is nonnegative for every \(t\), so its
    discriminant is nonpositive:
    \[
        4\expectation[XY]^2
        -
        4\expectation[X^2]\expectation[Y^2]
        \leq0.
    \]
    Taking square roots gives the inequality.
\end{snippetproof}

\section{Variance and covariance}

\begin{snippetdefinition}{variance-covariance-definition}{Variance and covariance}
    Let \((\Omega,\powerset(\Omega),\mathbb P)\) be a
    \countableprobabilityspace, and let
    \(X,Y\colon\Omega\fromto\realnumbers\) be real-valued
    \discreterandomvariable[discrete random variables] with finite second
    moments. The \emph{variance} of \(X\) is
    \[
        \variance(X)
        =
        \expectation[(X-\expectation[X])^2]
        =
        \expectation[X^2]-\expectation[X]^2.
    \]
    The \emph{covariance} of \(X\) and \(Y\) is
    \[
        \covariance(X,Y)
        =
        \expectation[(X-\expectation[X])(Y-\expectation[Y])]
        =
        \expectation[XY]-\expectation[X]\expectation[Y].
    \]
\end{snippetdefinition}

\begin{snippetproposition}{variance-covariance-density-formulas}{Density formulas for variance and covariance}
    Let \((\Omega,\powerset(\Omega),\mathbb P)\) be a
    \countableprobabilityspace, and let
    \(X,Y\colon\Omega\fromto\realnumbers\) be real-valued
    \discreterandomvariable[discrete random variables] with finite second
    moments. Then
    \[
        \variance(X)
        =
        \sum_{x\in\realnumbers}
        (x-\expectation[X])^2p_X(x),
    \]
    and
    \[
        \covariance(X,Y)
        =
        \sum_{(x,y)\in\realnumbers^2}
        (x-\expectation[X])(y-\expectation[Y])p_{X,Y}(x,y).
    \]
\end{snippetproposition}

\begin{snippetproof}{variance-covariance-density-formulas-proof}{variance-covariance-density-formulas}{Density formulas for variance and covariance}
    Apply the change-of-variable formula for the \expectedvalue to
    \(g(x)=(x-\expectation[X])^2\) and to
    \(h(x,y)=(x-\expectation[X])(y-\expectation[Y])\).
\end{snippetproof}

\begin{snippetproposition}{zero-variance-characterization}{Zero variance}
    Let \((\Omega,\powerset(\Omega),\mathbb P)\) be a
    \countableprobabilityspace, and let
    \(X\colon\Omega\fromto\realnumbers\) be a real-valued
    \discreterandomvariable with finite second moment. Then
    \[
        \variance(X)=0
        \quad\Longleftrightarrow\quad
        X\asequal \expectation[X].
    \]
\end{snippetproposition}

\begin{snippetproof}{zero-variance-characterization-proof}{zero-variance-characterization}{Zero variance}
    If \(X\asequal \expectation[X]\), then the defining formula gives
    \(\variance(X)=0\). Conversely, if \(\variance(X)=0\), then the
    nonnegative random variable \((X-\expectation[X])^2\) has expected value
    \(0\). By the basic property of the \expectedvalue, it is \(0\) almost
    surely, so \(X\asequal \expectation[X]\).
\end{snippetproof}

\begin{snippetproposition}{expected-value-minimizes-mean-square-error}{Expected value minimizes mean squared error}
    Let \((\Omega,\powerset(\Omega),\mathbb P)\) be a
    \countableprobabilityspace, and let
    \(X\colon\Omega\fromto\realnumbers\) be a real-valued
    \discreterandomvariable with finite second moment. Then, for every
    \(c\in\realnumbers\),
    \[
        \variance(X)\leq\expectation[(X-c)^2],
    \]
    with equality if and only if \(c=\expectation[X]\).
\end{snippetproposition}

\begin{snippetproof}{expected-value-minimizes-mean-square-error-proof}{expected-value-minimizes-mean-square-error}{Expected value minimizes mean squared error}
    Expanding around \(\expectation[X]\),
    \begin{align*}
        \expectation[(X-c)^2]
        &=
        \expectation[(X-\expectation[X]+\expectation[X]-c)^2]
        \\
        &=
        \variance(X)
        +
        2(\expectation[X]-c)\expectation[X-\expectation[X]]
        +
        (\expectation[X]-c)^2
        \\
        &=
        \variance(X)+(\expectation[X]-c)^2.
    \end{align*}
    The result follows.
\end{snippetproof}

\begin{snippetproposition}{variance-covariance-algebraic-properties}{Algebraic properties of variance and covariance}
    Let \((\Omega,\powerset(\Omega),\mathbb P)\) be a
    \countableprobabilityspace, and let
    \(X_1,\ldots,X_n\colon\Omega\fromto\realnumbers\) be real-valued
    \discreterandomvariable[discrete random variables] with finite second
    moments. Then:
    \begin{enumerate}
        \item for \(a,b\in\realnumbers\),
            \[
                \variance(aX_1+b)=a^2\variance(X_1);
            \]
        \item \(\variance(X_1)=\covariance(X_1,X_1)\);
        \item \(\covariance(X_1,X_2)=\covariance(X_2,X_1)\);
        \item for \(a,b\in\realnumbers\),
            \[
                \covariance(aX_1+bX_2,X_3)
                =
                a\covariance(X_1,X_3)+b\covariance(X_2,X_3);
            \]
        \item
            \[
                \variance\left(\sum_{k=1}^nX_k\right)
                =
                \sum_{k=1}^n\variance(X_k)
                +
                2\sum_{1\leq i<j\leq n}\covariance(X_i,X_j).
            \]
    \end{enumerate}
\end{snippetproposition}

\begin{snippetproof}{variance-covariance-algebraic-properties-proof}{variance-covariance-algebraic-properties}{Algebraic properties of variance and covariance}
    Each identity follows by expanding the definitions and using linearity of
    the \expectedvalue.
\end{snippetproof}

\begin{snippetproposition}{independence-implies-zero-covariance}{Independence implies zero covariance}
    Let \((\Omega,\powerset(\Omega),\mathbb P)\) be a
    \countableprobabilityspace, and let
    \(X,Y\colon\Omega\fromto\realnumbers\) be independent real-valued
    \discreterandomvariable[discrete random variables]. If
    \[
        \expectation[|X|]<+\infty,
        \qquad
        \expectation[|Y|]<+\infty,
    \]
    then \(X,Y\) have a \mixedmoment,
    \[
        \expectation[XY]=\expectation[X]\expectation[Y],
    \]
    and, if they have finite second moments,
    \[
        \covariance(X,Y)=0.
    \]
\end{snippetproposition}

\begin{snippetproof}{independence-implies-zero-covariance-proof}{independence-implies-zero-covariance}{Independence implies zero covariance}
    By independence, \(p_{X,Y}(x,y)=p_X(x)p_Y(y)\). Hence
    \[
        \expectation[|XY|]
        =
        \sum_{(x,y)\in\realnumbers^2}|xy|p_X(x)p_Y(y)
        =
        \left(\sum_{x\in\realnumbers}|x|p_X(x)\right)
        \left(\sum_{y\in\realnumbers}|y|p_Y(y)\right)
        <+\infty.
    \]
    The same factorization without absolute values gives
    \[
        \expectation[XY]
        =
        \expectation[X]\expectation[Y].
    \]
    Substituting this identity in the definition of covariance gives
    \(\covariance(X,Y)=0\).
\end{snippetproof}

\begin{snippetcorollary}{variance-sum-independent-random-variables}{Variance of a sum of independent random variables}
    Let \((\Omega,\powerset(\Omega),\mathbb P)\) be a
    \countableprobabilityspace, and let
    \(X_1,\ldots,X_n\colon\Omega\fromto\realnumbers\) be independent
    real-valued \discreterandomvariable[discrete random variables] with
    finite second moments. Then
    \[
        \variance\left(\sum_{k=1}^nX_k\right)
        =
        \sum_{k=1}^n\variance(X_k).
    \]
\end{snippetcorollary}

\begin{snippetproof}{variance-sum-independent-random-variables-proof}{variance-sum-independent-random-variables}{Variance of a sum of independent random variables}
    Distinct variables in the family have covariance \(0\). The result follows
    from the variance formula for a sum.
\end{snippetproof}

\begin{snippetexample}{zero-covariance-dependent-random-variables-example}{Zero covariance does not imply independence}
    Let \(X\) be a real-valued \discreterandomvariable with
    \[
        \mathbb P(X=-1)=\mathbb P(X=0)=\mathbb P(X=1)=\frac13,
    \]
    and let \(Y=X^2\). Then \(X\) and \(Y\) are not independent, because
    \(Y\) is determined by \(X\). However,
    \[
        \expectation[X]=0,
        \qquad
        \expectation[Y]=\expectation[X^2]=\frac23,
        \qquad
        \expectation[XY]=\expectation[X^3]=0.
    \]
    Therefore \(\covariance(X,Y)=0\).
\end{snippetexample}

\section{Correlation}

\begin{snippetproposition}{covariance-cauchy-schwarz-bound}{Cauchy-Schwarz bound for covariance}
    Let \((\Omega,\powerset(\Omega),\mathbb P)\) be a
    \countableprobabilityspace, and let
    \(X,Y\colon\Omega\fromto\realnumbers\) be real-valued
    \discreterandomvariable[discrete random variables] with finite second
    moments. Then
    \[
        |\covariance(X,Y)|
        \leq
        \sqrt{\variance(X)\variance(Y)}.
    \]
    If \(\variance(X)>0\) and \(\variance(Y)>0\), equality holds if and only
    if there exist \(a,b\in\realnumbers\), with \(a\neq0\), such that
    \[
        Y\asequal aX+b.
    \]
\end{snippetproposition}

\begin{snippetproof}{covariance-cauchy-schwarz-bound-proof}{covariance-cauchy-schwarz-bound}{Cauchy-Schwarz bound for covariance}
    Apply Cauchy-Schwarz to
    \[
        \overline X=X-\expectation[X],
        \qquad
        \overline Y=Y-\expectation[Y].
    \]
    This gives
    \[
        |\covariance(X,Y)|
        =
        |\expectation[\overline X\overline Y]|
        \leq
        \sqrt{\expectation[\overline X^2]\expectation[\overline Y^2]}
        =
        \sqrt{\variance(X)\variance(Y)}.
    \]
    Equality in Cauchy-Schwarz holds exactly when
    \(\overline Y\asequal a\overline X\) for some nonzero
    \(a\in\realnumbers\), which is equivalent to
    \(Y\asequal aX+b\) for
    \(b=\expectation[Y]-a\expectation[X]\).
\end{snippetproof}

\begin{snippetdefinition}{correlation-coefficient-definition}{Correlation coefficient}
    Let \((\Omega,\powerset(\Omega),\mathbb P)\) be a
    \countableprobabilityspace, and let
    \(X,Y\colon\Omega\fromto\realnumbers\) be real-valued
    \discreterandomvariable[discrete random variables] with finite second
    moments and positive variances. The \emph{correlation coefficient}
    between \(X\) and \(Y\) is
    \[
        \correlationcoefficient(X,Y)
        =
        \frac{\covariance(X,Y)}
        {\sqrt{\variance(X)\variance(Y)}}.
    \]
\end{snippetdefinition}

\begin{snippetproposition}{correlation-coefficient-range-and-equality}{Range and equality case for correlation}
    Let \((\Omega,\powerset(\Omega),\mathbb P)\) be a
    \countableprobabilityspace, and let
    \(X,Y\colon\Omega\fromto\realnumbers\) be real-valued
    \discreterandomvariable[discrete random variables] with finite second
    moments and positive variances. Then
    \[
        -1\leq\correlationcoefficient(X,Y)\leq1.
    \]
    Moreover, \(\correlationcoefficient(X,Y)=1\) if and only if
    \(Y\asequal aX+b\) for some \(a>0\) and \(b\in\realnumbers\), while
    \(\correlationcoefficient(X,Y)=-1\) if and only if
    \(Y\asequal aX+b\) for some \(a<0\) and \(b\in\realnumbers\).
\end{snippetproposition}

\begin{snippetproof}{correlation-coefficient-range-and-equality-proof}{correlation-coefficient-range-and-equality}{Range and equality case for correlation}
    The range follows by dividing the covariance bound by
    \(\sqrt{\variance(X)\variance(Y)}\). The equality cases are the equality
    cases of the covariance bound, with the sign determined by the sign of
    the covariance.
\end{snippetproof}

\end{document}
